\clearpage
\section{Paired Behavioral Opportunity Across Model Configurations}
\label{app:paired-opportunity}

\subsection{A shared, metadata-matched patient cohort}

The paired-opportunity study evaluates 100 fresh test patients, each with
one Consolidation-positive and one Consolidation-negative study. Patients
are absent from the entire original probe manifest and all preceding
intervention cohorts, including the earlier 50-pair closure study.
Within each patient, the two studies match on view, sex, integer age, and
all 13 other source disease labels. The fixed selection retains complete
source records and patient/image order. Matching defines the paired
population; unrecorded longitudinal differences can remain.

Figure~\ref{fig:paired-prompt} gives the common clinical prompt. Qwen is
evaluated first, followed by a registered LLaVA comparison on exactly the
same 200 images. The model-specific processor and chat template wrap that
prompt. Qwen retains its accepted image processing; LLaVA uses its snapshot
processor's RGB conversion, 336-pixel shortest-edge resize, $336\times336$
center crop, CLIP normalization, and default visual feature selection.

\begin{figure}[H]
\centering
\begin{minipage}{0.95\linewidth}
\small
\fbox{\parbox{0.95\linewidth}{Is there consolidation in this chest radiograph? Answer yes or no.}}
\end{minipage}
\caption{\textbf{Both model configurations evaluate the same paired-input question.} The clinical prompt is fixed across the positive and negative studies and across models.}
\label{fig:paired-prompt}
\end{figure}

Each model supplies its first-answer yes-minus-no margin, taking the
maximum over accepted singleton case and leading-space variants for each
answer. The paired gap subtracts the negative study's margin from the
positive study's margin in that model's own coordinates. The sigmoid of
each image's margin gives its yes probability; the probability gap
subtracts the two probabilities. This is a clean-response measurement,
with no representation intervention. The Qwen and LLaVA runs each contain
200 image outcomes and take 31 and 38 seconds, respectively, on one A100.

\subsection{Model-specific opportunity and the direct paired comparison}

The registered opportunity endpoint is the patient-weighted mean margin
gap within each model. Its one-sided 95\% lower bound is the fifth
percentile of 5,000 patient-bootstrap estimates. A strictly positive bound
is required for positive opportunity. Both analyses use the identical
stored resampling indices, generated with seed 20260912, and retain all
100 patients irrespective of the sign of their gap.
Table~\ref{tab:paired-gap-intervals} presents these model-specific endpoints
and their descriptive two-sided percentile intervals.

\begin{table}[H]
\centering
\caption{Uncertainty in paired Consolidation response.}
\label{tab:paired-gap-intervals}
\resizebox{\ifdim\width>\linewidth\linewidth\else\width\fi}{!}{%
\begin{tabular}{lrrr}
\toprule
Model & Mean margin gap & One-sided 95\% lower bound & Descriptive 95\% interval \\
\midrule
Qwen2.5-VL-7B & $-0.1175$ & $-0.3987$ & $[-0.4500,0.2088]$ \\
LLaVA-1.5-7B & $-0.0712$ & $-0.1338$ & $[-0.1463,0.0012]$ \\
\bottomrule
\end{tabular}%
}
\end{table}

Both lower bounds are negative, so neither configuration meets the
opportunity criterion. The two-sided intervals include positive and
negative mean responses: positive opportunity remains unresolved in this
metadata-defined population. The model-specific margin scales are not
subtracted from one another to infer a model difference.

For that direct comparison, within-patient concordance assigns one to a
positive gap, zero to a negative gap, and one half to a tie. A tie is exact
equality of the positive and negative margins under the registered scoring
configured for BF16. A model's concordance is the mean of these patient
scores. The cross-model endpoint
first subtracts Qwen's score from LLaVA's for each patient, then averages
over the shared cohort. Its descriptive 95\% interval reuses the same
5,000 patient-bootstrap draws. Table~\ref{tab:paired-concordance} shows the
gap-sign counts and the direct contrast.

\begin{table}[H]
\centering
\caption{Within-patient concordance and its direct model contrast.}
\label{tab:paired-concordance}
\resizebox{\ifdim\width>\linewidth\linewidth\else\width\fi}{!}{%
\begin{tabular}{lrrrrr}
\toprule
Model or contrast & Positive & Negative & Tied & Concordance & 95\% interval \\
\midrule
Qwen2.5-VL-7B & 54 & 45 & 1 & 0.545 & --- \\
LLaVA-1.5-7B & 36 & 45 & 19 & 0.455 & --- \\
LLaVA $-$ Qwen & --- & --- & --- & $-0.090$ & $[-0.210,0.025]$ \\
\bottomrule
\end{tabular}%
}
\end{table}

The direct interval spans zero, leaving the concordance difference
unresolved. The sign counts account for every selected patient, and ties
receive the same half-credit treatment in both models. The direct contrast
uses the paired patient scores, rather than comparing the significance of
the separate model-specific opportunity tests.

\subsection{Probability gaps and source-label strata}

Table~\ref{tab:paired-strata} reports complementary descriptive summaries.
The cohort contains 52 pairs whose negative image is labelled
\texttt{No Finding} and 48 whose negative image has other source findings.
Stratum means use the within-model logit gap; the full-cohort probability
gap uses the difference of sigmoid scores defined above.

\begin{table}[H]
\centering
\caption{Paired response across source-label strata.}
\label{tab:paired-strata}
\resizebox{\ifdim\width>\linewidth\linewidth\else\width\fi}{!}{%
\begin{tabular}{lrrr}
\toprule
Model & \shortstack{Full-cohort\\probability gap} & \shortstack{No Finding\\margin gap ($n=52$)} & \shortstack{Other findings\\margin gap ($n=48$)} \\
\midrule
Qwen2.5-VL-7B & $-0.01743$ & $-0.21154$ & $-0.01563$ \\
LLaVA-1.5-7B & $-0.01508$ & $-0.15385$ & $0.01823$ \\
\bottomrule
\end{tabular}%
}
\end{table}

Mean probability gaps are negative in both configurations. The stratum
summaries describe the selected population and leave the registered
full-cohort opportunity decisions unchanged. Taken together, the mean-gap
uncertainty and the direct concordance comparison identify the behavioral
opportunity that remains unresolved before interpreting component-based
recovery. Because this cohort differs from the original 50-pair closure
cohort, it contextualizes that study rather than identifying the cause of
its closure result.

\FloatBarrier
